\section{Comparación crítica con otra plataforma de video}

Para evitar comparar solo tecnologías, contrastamos el caso de Netflix con la
infraestructura de procesamiento de video que Google desplegó para cargas como
YouTube. Ambos sistemas trabajan con video a gran escala, pero la evidencia
pública describe énfasis diferentes.

\begin{center}
\small
\begin{tabularx}{\linewidth}{@{}p{2.5cm}XX@{}}
\toprule
\textbf{Criterio} & \textbf{Netflix} & \textbf{YouTube/Google} \\
\midrule
Problema observado & Servicios y datos para producir, descubrir y entregar
contenido. & Procesamiento y transcodificación de video a escala de centro de
datos. \\
Arquitectura pública & Microservicios, cachés y almacenes distribuidos; gRPC
predomina en comunicación backend. & Software distribuido coordinado con
aceleradores VCU especializados. \\
Comunicación & Contratos gRPC/Protocol Buffers en los casos documentados; otras
interfaces también existen. & La publicación consultada no permite afirmar un
protocolo RPC concreto. \\
Disponibilidad & EVCache mantiene operaciones locales aunque la réplica remota
se atrase. & Escalado horizontal y adaptación de la infraestructura ante fallos
de procesamiento. \\
Trade-off & Baja latencia y disponibilidad para ciertos datos a cambio de
consistencia eventual. & Mayor eficiencia de transcodificación a cambio de
hardware y co-diseño especializados. \\
\bottomrule
\end{tabularx}
\end{center}

Google reportó mejoras de eficiencia de 20 a 33 veces frente a su plataforma
anterior al desplegar aceleradores de video a escala. No concluimos que esa
arquitectura sea universalmente superior: resuelve una parte distinta del
problema y exige especialización de hardware \cite{google-vcu}.

\paragraph{RPC, gRPC y RMI.} RPC es el modelo general de invocación remota;
gRPC es un framework multilenguaje basado normalmente en contratos protobuf;
Java RMI permite invocar objetos remotos entre programas Java. RMI podría ser
razonable en un sistema homogéneo Java y controlado, pero no debe atribuirse a
Netflix ni a YouTube sin evidencia.
